% !TeX root = ../main.tex
\chapter{Methodology}
\label{chap:methodology}

This chapter formalizes the econometric and machine learning framework for multi-horizon soybean yield anomaly forecasting across the Midwestern Corn Belt. Building on the 75-year panel characterized in Chapter~\ref{chap:data_and_study_area}, we formalize the predictive pipeline without anticipating empirical findings. The framework pursues two objectives: determining the operational lead time where meteorological predictors yield significant skill over secular trend baselines, and evaluating whether non-linear algorithmic complexity improves performance under strict out-of-sample expanding-window validation.

% ==============================================================================
% SECTION 4.1: PROBLEM FORMULATION AND MULTI-HORIZON FORECAST SETUP
% ==============================================================================

\section{Problem Formulation and Multi-Horizon Forecast Setup}
\label{sec:meth_problem_setup}

\subsection{Continuous Target Formulation and Train-Only Detrending}
\label{sec:meth_target_formulation}

Realized county soybean yields decompose additively into a secular technological trend and an annual climatic anomaly (Section~\ref{sec:data_yield_decomposition}):
\begin{equation}
    y_{c,Y} = \tau_{c,Y} + \epsilon_{c,Y}
    \label{eq:yield_additive_decomp}
\end{equation}
where $c \in \{1, \dots, C\}$ is the county index ($C = 135$), $Y \in \{1951, \dots, 2025\}$ denotes harvest crop year, $\tau_{c,Y}$ designates the secular technological trend, and $\epsilon_{c,Y}$ isolates the weather-induced anomaly. Predictive models target the continuous anomaly $\epsilon_{c,Y} \in \mathbb{R}$ directly.

Detrending parameters are estimated strictly on training observations to prevent look-ahead leakage. For harvest year $Y$, county-specific linear trend parameters $(\hat{\alpha}_c^{(Y)}, \hat{\beta}_c^{(Y)})$ are fitted over historical campaigns $\mathcal{T}_{\text{train}}(Y) = \{1951, \dots, Y-1\}$:
\begin{equation}
    \hat{\tau}_{c,Y} = \hat{\alpha}_{c}^{(Y)} + \hat{\beta}_{c}^{(Y)} Y, \quad (\hat{\alpha}_{c}^{(Y)}, \hat{\beta}_{c}^{(Y)}) = \arg\min_{\alpha, \beta} \sum_{t \in \mathcal{T}_{\text{train}}(Y)} \left( y_{c,t} - (\alpha + \beta t) \right)^2
    \label{eq:train_only_trend}
\end{equation}
Training targets are $\epsilon_{c,t} = y_{c,t} - \hat{\tau}_{c,t}$. On held-out evaluation year $Y$, models predict continuous anomaly $\hat{\epsilon}_{c,Y}$, and total yield levels are reconstructed post-hoc:
\begin{equation}
    \hat{y}_{c,Y} = \hat{\tau}_{c,Y} + \hat{\epsilon}_{c,Y}
    \label{eq:yield_reconstruction}
\end{equation}

\subsection{Temporal Resolution Mismatch: Bridging Daily Weather to Annual Scalar Yield}
\label{sec:meth_temporal_mismatch}

The prediction target is a single annual scalar per county ($\epsilon_{c,Y} \in \mathbb{R}$), whereas meteorological forcing operates as daily time series across the crop year ($K = 16$ indicators, Section~\ref{sec:meth_meteo_pool}). Ingesting raw daily weather across 365 days generates $365 \times 16 \approx 5{,}840$ features per county-year. Against a baseline training pool of $N_{\text{train}} = 6{,}075$ county-years (1951--1995), this unconstrained parameterization causes severe collinearity and out-of-sample variance inflation.

To preserve phenological timing while restricting dimensionality, daily weather fields are aggregated into structured monthly blocks:
\begin{equation}
    w_{k, j} = 
    \begin{cases}
        \displaystyle \sum_{d \in \mathcal{D}_j} v_{k, d}, & \text{for cumulative flux indicators } k \in \mathcal{K}_{\text{flux}} \\[2.5ex]
        \displaystyle \frac{1}{|\mathcal{D}_j|} \sum_{d \in \mathcal{D}_j} v_{k, d}, & \text{for environmental state indicators } k \in \mathcal{K}_{\text{state}}
    \end{cases}
    \label{eq:monthly_aggregation_rule}
\end{equation}
where $\mathcal{K}_{\text{flux}} = \{P, SSRD, ET_0, P-ET_0, GDD, HD_{30}, HD_{35}\}$ denotes flux integrals, $\mathcal{K}_{\text{state}} = \{T_{\text{mean}}, T_{\text{max}}, T_{\text{min}}, T_{\text{dew}}, SM_1, SM_2, SM_3, VPD, SM_{\text{root}}\}$ designates environmental states, and $\mathcal{D}_j$ denotes calendar days in campaign month $j \in \{1, \dots, 12\}$. For each elapsed month $j$, this defines a 16-dimensional meteorological summary vector $\mathbf{w}_{c,Y,j} \in \mathbb{R}^{16}$.

\subsection{Crop Campaign Timeline and Multi-Horizon Feature Expansion}
\label{sec:meth_timeline_horizons}

The crop campaign spans November 1 (year $Y-1$) to October 31 (year $Y$). Forecasts are generated at the end of each campaign month $m \in \{1, \dots, 12\}$, establishing lead times $H = 13 - m$ prior to harvest.

At forecast horizon $H$, tabular feature vector $\mathbf{x}_{c,Y}^{(H)}$ concatenates normalized spatial coordinates $(\tilde{\text{Lat}}_c, \tilde{\text{Lon}}_c)$ (Section~\ref{sec:meth_spatial_covariates}) with elapsed monthly weather summaries:
\begin{equation}
    \mathbf{x}_{c,Y}^{(H)} = \left[ \, \tilde{\text{Lat}}_c, \, \tilde{\text{Lon}}_c, \, \mathbf{w}_{c,Y, 1}^T, \, \dots, \, \mathbf{w}_{c,Y, m}^T \, \right]^T \in \mathbb{R}^{P_H}
    \label{eq:tabular_feature_vector}
\end{equation}
Feature dimensionality expands as meteorological records accumulate:
\begin{equation}
    P_H = 
    \begin{cases}
        P_{\text{geo}} = 2, & \text{for } H = 12 \quad (m = 1, \text{ zero weather}) \\[1ex]
        P_{\text{geo}} + K \cdot m = 2 + 16 \cdot m, & \text{for } H \in \{11, \dots, 1\} \quad (m = 2, \dots, 12)
    \end{cases}
    \label{eq:feature_dimension_formula}
\end{equation}
This schedule structures information accumulation:
\begin{itemize}
    \item \textbf{Horizon $H = 12$ ($m = 1$, Nov 30): Zero-Weather Benchmark.} Ingests spatial coordinates only ($P_{12} = 2, P_{\text{met}} = 0$). Candidate models output $\hat{\epsilon}_{c,Y}^{(12)} = 0$, defining the unaugmented secular trend baseline ($\hat{y}_{c,Y}^{(12)} = \hat{\tau}_{c,Y}$, $R^2_{\text{OOS}} = 0$).
    \item \textbf{Horizons $H \in \{11, \dots, 1\}$ ($m \in \{2, \dots, 12\}$): Cumulative Weather Ingestion.} Models ingest observed weather through Month $m$, expanding feature dimensionality from $P_{11} = 34$ to $P_1 = 194$ in October.
\end{itemize}
This formulation evaluates whether autumn recharge ($H = 11$) provides predictive skill over trend inertia, while identifying the lead time where weather forcing dominates. Dimensionality remains bounded by $P_1 = 194$.

\subsection{Direct Horizon-Specific Modeling Paradigm}
\label{sec:meth_direct_horizon}

We implement direct multi-horizon forecasting. For each lead time $H \in \{12, \dots, 1\}$, an independent model $\mathcal{M}_H$ is trained on horizon-specific feature vector $\mathbf{x}_{c,Y}^{(H)}$:
\begin{equation}
    \hat{\epsilon}_{c,Y}^{(H)} = \mathcal{M}_H\left(\mathbf{x}_{c,Y}^{(H)}; \, \mathbf{\Theta}_H\right)
    \label{eq:direct_model_instance}
\end{equation}
where $\mathbf{\Theta}_H$ denotes the parameter vector. Direct forecasting prevents compounding error propagation from iterated rollouts. It allows model weights to reconfigure across developmental stages, prioritizing early soil moisture at long lead times and reproductive heat stress during late-summer origins.

\subsection{Pooled Panel Framework and Disaggregated County Evaluation}
\label{sec:meth_pooled_panel}

Models are estimated within a pooled cross-sectional panel framework across all $C = 135$ counties simultaneously. Separate county-by-county estimation with 45 to 74 historical observations causes parameter collapse ($P_H > T$) against feature vectors with up to 194 variables. Pooling expands sample size to $N_{\text{train}} = 135 \times |\mathcal{T}_{\text{train}}(Y)|$ observations ($N_{\text{train}} = 6{,}075$ on the baseline 1951--1995 partition, expanding to $8{,}100$ at fold 15 in 2025). 

Importantly, while estimation parameters are pooled, predictions remain county-specific: for each county $c$ and year $Y$, the model generates an individual anomaly forecast $\hat{\epsilon}_{c,Y}^{(H)}$ by evaluating localized weather series and spatial coordinates $(\tilde{\text{Lat}}_c, \tilde{\text{Lon}}_c)$. County yield forecasts are reconstructed using county-specific trends $\hat{\tau}_{c,Y}$. To avoid reporting 135 fragmented performance scorecards, empirical skill is evaluated both at the pooled panel level and via annual regional averages (Section~\ref{sec:meth_evaluation}).

% ==============================================================================
% SECTION 4.2: SPATIAL COVARIATES AND INPUT REPRESENTATION
% ==============================================================================

\section{Spatial Covariates and Input Representation}
\label{sec:meth_input_representation}

\subsection{Continuous Spatial Coordinates versus Administrative State Dummies}
\label{sec:meth_spatial_covariates}

Spatial location is encoded exclusively through normalized continuous centroid coordinates:
\begin{equation}
    \mathbf{s}_c = (\text{Lat}_c, \, \text{Lon}_c) \in \mathbb{R}^2, \quad P_{\text{geo}} = 2
    \label{eq:spatial_coordinates}
\end{equation}
Coordinates are scaled to $[0, 1]$ using training bounds:
\begin{equation}
    \tilde{\text{Lat}}_c = \frac{\text{Lat}_c - \min_{i} \text{Lat}_i}{\max_{i} \text{Lat}_i - \min_{i} \text{Lat}_i}, \quad
    \tilde{\text{Lon}}_c = \frac{\text{Lon}_c - \min_{i} \text{Lon}_i}{\max_{i} \text{Lon}_i - \min_{i} \text{Lon}_i}
    \label{eq:coord_normalization}
\end{equation}

Administrative state dummies are excluded. Political state borders do not reflect agro-ecological boundaries. Contiguous counties across state lines share identical soils, microclimates, and topography, whereas counties within the same state span substantial latitudinal gradients. Categorical state dummies introduce artificial discontinuities at administrative borders. Continuous coordinates $(\tilde{\text{Lat}}_c, \tilde{\text{Lon}}_c)$ allow non-linear estimators to learn smooth biophysical response surfaces, capturing spatial autocorrelation and photoperiod gradients natively.

\subsection{Agrometeorological Predictor Space and Functional Partitioning}
\label{sec:meth_meteo_pool}

The meteorological predictor space builds upon the $K = 16$ daily environmental indicators derived from ECMWF ERA5-Land reanalysis in Section~\ref{sec:data_derived_weather} and cataloged in Table~\ref{tab:weather_variables}. For empirical estimation, we formalize their functional partitioning:
\begin{enumerate}
    \item \textbf{Flux Integrals ($\mathcal{K}_{\text{flux}}$, 7 variables):} Total precipitation ($P$), surface solar radiation downwards ($SSRD$), reference evapotranspiration ($ET_0$), climatic water balance ($P - ET_0$), growing degree days ($GDD$), moderate heat days ($HD_{30}$), and extreme heat days ($HD_{35}$). These indicators represent volumetric and energy accumulation, aggregated via cumulative summation over temporal windows (Eq.~\eqref{eq:monthly_aggregation_rule}).
    \item \textbf{Environmental State Variables ($\mathcal{K}_{\text{state}}$, 9 variables):} Daily mean, maximum, minimum, and dewpoint temperatures ($T_{\text{mean}}, T_{\text{max}}, T_{\text{min}}, T_{\text{dew}}$), volumetric soil water across three horizons ($SM_1, SM_2, SM_3$), vapor pressure deficit ($VPD$), and root-zone soil moisture ($SM_{\text{root}}$). These indicators reflect continuous background ambient stress, aggregated via arithmetic time-averaging (Eq.~\eqref{eq:monthly_aggregation_rule}).
\end{enumerate}
This functional grouping maps daily physical reanalysis directly into structured feature matrices without narrative duplication across chapters.

\subsection{Methodological Trade-offs in Temporal Aggregation Granularity}
\label{sec:meth_aggregation_tradeoffs}

\citet{Xie2025} and \citet{Yinetal2026} establish that sub-monthly pentad (5-day) and decad (10-day) aggregations capture short-duration heat waves and heavy precipitation events that are smoothed within monthly means. However, in tabular modeling architectures where temporal blocks are concatenated horizontally across columns (Eq.~\eqref{eq:tabular_feature_vector}), sub-monthly partitions induce rapid feature explosion: a 73-pentad schedule yields $2 + 16 \times 73 = 1{,}170$ features, inducing severe multicollinearity and variance inflation on historical sample sizes \citep{Sharmaetal2025, Hoffman2020}. Monthly blocks maintain dimensional stability ($P_H \le 194$). To evaluate high-frequency sequential dynamics without expanding column dimensions, we deploy sequential deep architectures within the candidate portfolio.

% ==============================================================================
% SECTION 4.3: CANDIDATE FORECASTING MODELS AND COMPLEXITY LADDER
% ==============================================================================

\section{Candidate Forecasting Models and Complexity Ladder}
\label{sec:meth_models}

To isolate the predictive value of meteorological features and test the utility of algorithmic complexity, we structure candidate algorithms into a four-tier complexity ladder.

\subsection{Reference Baseline: Naive Deterministic Trend Benchmark}
\label{sec:meth_baselines}

To establish the statistical boundary where meteorological predictors provide genuine predictive skill, we formulate the naive deterministic trend benchmark ($\mathcal{B}_{\text{naive}}$):
\begin{equation}
    \hat{\epsilon}_{c,Y} = 0 \iff \hat{y}_{c,Y} = \hat{\tau}_{c,Y}
    \label{eq:naive_baseline}
\end{equation}
This reference benchmark assumes zero weather-induced anomaly, projecting future yields strictly along the county's historical technological trend. At Month 1 ($H=12$), where zero meteorological predictors are ingested ($P_{\text{met}} = 0$), this benchmark defines the reference origin with an out-of-sample skill score $R^2_{\text{OOS}} = 0$. Any positive $R^2_{\text{OOS}} > 0$ achieved at subsequent horizons ($H \le 11$) demonstrates genuine value from environmental forcing beyond historical inertia.

\subsection{Regularized Linear Model: ElasticNet}
\label{sec:meth_elasticnet}

ElasticNet minimizes a penalized sum of squared residuals combining $L_1$ (Lasso) and $L_2$ (Ridge) penalties:
\begin{equation}
    \min_{\mathbf{\beta}, \beta_0} \frac{1}{2 N} \sum_{i=1}^{N} \left( \epsilon_i - (\beta_0 + \mathbf{x}_i^T \mathbf{\beta}) \right)^2 + \lambda \left( \alpha \|\mathbf{\beta}\|_1 + \frac{1 - \alpha}{2} \|\mathbf{\beta}\|_2^2 \right)
    \label{eq:elasticnet_objective}
\end{equation}
where $\lambda > 0$ controls shrinkage intensity and $\alpha \in [0, 1]$ governs penalty mixing. ElasticNet performs automatic feature selection via $L_1$ sparsity while grouping collinear agrometeorological predictors stably through the $L_2$ quadratic term.


\subsection{Decision Tree Ensembles: Random Forest and XGBoost}
\label{sec:meth_trees}

We specify two complementary tree ensembles:
\begin{enumerate}
    \item \textbf{Random Forest (Bagging):} Averages $B$ deep regression trees trained on bootstrap samples \citep{Sharmaetal2025, Torsonietal2023}:
    \begin{equation}
        \hat{\epsilon}_{c,Y} = \frac{1}{B} \sum_{b=1}^B T_b\left(\mathbf{x}_{c,Y}; \, \mathbf{\Theta}_b\right)
        \label{eq:rf_prediction}
    \end{equation}
    At each candidate split, candidate partitioning minimizes node variance over a random feature subset $p_{\text{split}} \le P_H$. Random feature subsampling decorrelates individual trees, reducing ensemble variance on high-dimensional weather inputs.
    \item \textbf{XGBoost (Gradient Tree Boosting):} Fits an additive sequence of $M$ shallow regression trees minimizing a regularized loss \citep{Torsonietal2023}:
    \begin{equation}
        \mathcal{L}^{(m)} = \sum_{i=1}^{N} \left[ g_i f_m(\mathbf{x}_i) + \frac{1}{2} h_i f_m(\mathbf{x}_i)^2 \right] + \gamma_{\text{split}} T_{\text{leaves}} + \frac{1}{2} \lambda_{\text{reg}} \sum_{j=1}^{T_{\text{leaves}}} w_j^2 + \alpha_{\text{reg}} \sum_{j=1}^{T_{\text{leaves}}} |w_j|
        \label{eq:xgboost_objective}
    \end{equation}
    where $g_i$ and $h_i$ are first and second order loss gradients, $T_{\text{leaves}}$ is leaf count, and $w_j$ are leaf weights. Leaf complexity is constrained by split threshold $\gamma_{\text{split}}$, $L_2$ regularization $\lambda_{\text{reg}}$, and $L_1$ penalty $\alpha_{\text{reg}}$ \citep{Geron2022}. Column subsampling randomly restricts predictors at each tree level to prevent dominance by individual heat-stress indicators \citep{James2023}. Sequential updates employ shrinkage rate $\eta \in (0, 1]$:
    \begin{equation}
        \hat{\epsilon}_i^{(m)} = \hat{\epsilon}_i^{(m-1)} + \eta f_m(\mathbf{x}_i)
        \label{eq:xgboost_update}
    \end{equation}
\end{enumerate}

\subsection{Sequential Deep Learning: Long Short-Term Memory Network}
\label{sec:meth_lstm}

While tabular algorithms require flattened two-dimensional design matrices $\mathbf{x}_{c,Y}^{(H)} \in \mathbb{R}^{P_H}$ where temporal blocks expand horizontally across columns, sequential deep architectures decouple input feature dimensionality from temporal sequence length \citep{Khaki2020, Shooketal2021}. 

Instead of concatenating temporal blocks into expanding feature vectors, the Long Short-Term Memory (LSTM) network operates directly on a three-dimensional sequence tensor:
\begin{equation}
    \mathbf{X}^{(H)} \in \mathbb{R}^{N_{\text{train}} \times T_H \times D_{\text{in}}}
    \label{eq:tensor_dimensions}
\end{equation}
where $N_{\text{train}}$ is sample size, $T_H$ is sequence length across elapsed campaign steps at horizon $H$, and $D_{\text{in}} = 19$ represents the input channel dimensionality. Following the empirical spatial encoding evaluation in Section~\ref{sec:meth_input_representation}, $D_{\text{in}}$ comprises the $K = 17$ agrometeorological indicators (including acute heat stress $HD_{30}$ and $HD_{35}$) concatenated with the two normalized continuous spatial coordinates $(\tilde{\text{Lat}}_c, \tilde{\text{Lon}}_c)$ ingested at every sequence step. This enables recurrent memory cells to modulate hydrothermal responses conditionally on geographical latitude and longitude without requiring discrete county embeddings.

To resolve sub-monthly extreme events while maintaining computational tractability, meteorological series are structured into 10-day decad intervals ($W^* = 10\text{d}$). Across the 12-month campaign (November 1 to October 31), the full season spans 35 decad steps ($T_1 = 35$ at harvest $H=1$, scaling down to $T_{11} = 3$ at $H=11$). February is standardized to 28 days to preserve invariant sequence lengths across leap and non-leap years. This decad resolution captures short-duration heat waves and moisture deficits that are smoothed by monthly averages \citep{Yinetal2026}.

At each sequence step $t \in \{1, \dots, T_H\}$, the recurrent cell ingests the input vector $\mathbf{z}_t \in \mathbb{R}^{19}$ and updates its internal state via gated recurrent transitions:
\begin{align}
    \mathbf{f}_t &= \sigma\left(\mathbf{W}_f \mathbf{z}_t + \mathbf{U}_f \mathbf{h}_{t-1} + \mathbf{b}_f\right) \label{eq:lstm_f} \\
    \mathbf{i}_t &= \sigma\left(\mathbf{W}_i \mathbf{z}_t + \mathbf{U}_i \mathbf{h}_{t-1} + \mathbf{b}_i\right) \label{eq:lstm_i} \\
    \tilde{\mathbf{c}}_t &= \tanh\left(\mathbf{W}_c \mathbf{z}_t + \mathbf{U}_c \mathbf{h}_{t-1} + \mathbf{b}_c\right) \label{eq:lstm_c_tilde} \\
    \mathbf{c}_t &= \mathbf{f}_t \odot \mathbf{c}_{t-1} + \mathbf{i}_t \odot \tilde{\mathbf{c}}_t \label{eq:lstm_c} \\
    \mathbf{o}_t &= \sigma\left(\mathbf{W}_o \mathbf{z}_t + \mathbf{U}_o \mathbf{h}_{t-1} + \mathbf{b}_o\right) \label{eq:lstm_o} \\
    \mathbf{h}_t &= \mathbf{o}_t \odot \tanh\left(\mathbf{c}_t\right) \label{eq:lstm_h}
\end{align}
where $\mathbf{h}_t \in \mathbb{R}^{d_h}$ denotes the hidden state vector, $\mathbf{c}_t \in \mathbb{R}^{d_h}$ represents the internal cell state, and $\sigma(\cdot)$ is the sigmoid activation function. Forget gate $\mathbf{f}_t$ regulates historical information retention; input gate $\mathbf{i}_t$ modulates candidate cell updates $\tilde{\mathbf{c}}_t$; and output gate $\mathbf{o}_t$ filters the emitted hidden representation.

To synthesize information across all elapsed sequence steps without compressing the entire growing season into the single final state $\mathbf{h}_{T_H}$, we implement a Bahdanau temporal attention layer:
\begin{align}
    e_t &= \mathbf{v}_a^T \tanh\left(\mathbf{W}_a \mathbf{h}_t + \mathbf{b}_a\right) \label{eq:attention_score} \\
    \alpha_t &= \frac{\exp(e_t)}{\sum_{k=1}^{T_H} \exp(e_k)}, \quad \sum_{t=1}^{T_H} \alpha_t = 1 \label{eq:attention_weights} \\
    \mathbf{c}_{\text{ctx}} &= \sum_{t=1}^{T_H} \alpha_t \mathbf{h}_t \in \mathbb{R}^{d_h} \label{eq:context_vector}
\end{align}
where $\mathbf{W}_a \in \mathbb{R}^{d_a \times d_h}$ and $\mathbf{v}_a \in \mathbb{R}^{d_a}$ parameterize the attention projection ($d_a = 32$). Attention weights $\alpha_t$ quantify the relative contribution of each 10-day window, providing physiological transparency by mapping predictive focus across vegetative and reproductive growth stages.

We evaluate two projection head architectures mapping context vector $\mathbf{c}_{\text{ctx}}$ to scalar anomaly prediction $\hat{\epsilon}_{c,Y}^{(H)}$:
\begin{enumerate}
    \item \textbf{Non-Linear PReLU Head (Champion Specification):}
    \begin{equation}
        \hat{\epsilon}_{c,Y}^{(H)} = \mathbf{W}_2 \operatorname{PReLU}\left(\mathbf{W}_1 \left[ \mathbf{c}_{\text{ctx}}^T, \, \tilde{\text{Lat}}_c, \, \tilde{\text{Lon}}_c \right]^T + \mathbf{b}_1\right) + b_2
        \label{eq:lstm_prelu_head}
    \end{equation}
    where $\operatorname{PReLU}(z) = \max(0, z) + a \min(0, z)$ learns parameter $a$ during backpropagation. The parametric rectified linear unit captures asymmetric crop responses, allowing the network to apply steep penalizations to acute heat and moisture deficits while maintaining moderate slopes for benign climatic deviations.
    \item \textbf{Linear Projection Head (Parsimonious Benchmark):}
    \begin{equation}
        \hat{\epsilon}_{c,Y}^{(H)} = \mathbf{W}_{\text{out}} \left[ \mathbf{c}_{\text{ctx}}^T, \, \tilde{\text{Lat}}_c, \, \tilde{\text{Lon}}_c \right]^T + b_{\text{out}}
        \label{eq:lstm_linear_head}
    \end{equation}
    establishing the linear mapping baseline directly from the recurrent context vector.
\end{enumerate}

To evaluate the contribution of historical yield persistence, we formulate two parallel information tiers: \textit{Category B (Weather-Only)}, where inputs are strictly meteorological; and \textit{Category C (Integrated)}, where the historical detrended yield anomaly $\epsilon_{c, Y-1}$ (or $\epsilon_{c, Y-2}$ at early horizons $H \ge 8$, adhering to USDA NASS release timing) is concatenated directly to the dense projection layer $[\mathbf{c}_{\text{ctx}}^T, \tilde{\text{Lat}}_c, \tilde{\text{Lon}}_c, \epsilon_{c, \text{lag}}]^T$.

As a preliminary architectural ablation during validation, we evaluated the Gated Recurrent Unit (GRU) \citep{Geron2022}. Although the GRU offers a 25\% parameter reduction, its single hidden state without a dedicated cell state $\mathbf{c}_t$ proved insufficient for retaining cumulative soil moisture deficits across 35 decad steps, yielding inferior validation RMSE ($>4.91\text{ bu/ac}$ vs $4.66\text{ bu/ac}$ for LSTM). Consequently, GRU was discarded from blind out-of-sample test promotion.

Network parameters are optimized via AdamW with initial learning rate $\eta = 3 \times 10^{-4}$ and weight decay $10^{-4}$ \citep{Geron2022}. Learning rates decay dynamically via a plateau scheduler ($\text{factor} = 0.5$, $\text{patience} = 3$), and training terminates via early stopping ($\text{patience} = 7$) on train-only folds. Table~\ref{tab:model_summary} summarizes the complete forecasting portfolio.

\begin{table}[htbp]
\centering
\singlespacing
\footnotesize
\setlength{\tabcolsep}{3pt}
\renewcommand{\arraystretch}{1.12}
\caption{Comparative architectural summary of candidate forecasting models.}
\label{tab:model_summary}
\begin{tabularx}{\textwidth}{l p{2.2cm} p{2.0cm} p{2.4cm} p{2.4cm} X}
\toprule
\textbf{Model} & \textbf{Family} & \textbf{Input Space} & \textbf{Regularization} & \textbf{Tuning Parameters} & \textbf{Agronomic Function} \\
\midrule
$\mathcal{B}_{\text{naive}}$ & Secular Baseline & None ($P_{\text{met}}=0$) & Deterministic & None & Zero-weather trend baseline ($R^2_{\text{OOS}}=0$). \\
ElasticNet & Regularized Linear & Tabular 2D ($P_H$) & $L_1 + L_2$ convex loss & $\lambda$, $\alpha$ & Linear baseline with collinearity grouping. \\
Random Forest & Bagging Trees & Tabular 2D ($P_H$) & $p_{\text{split}}$, leaf size & $p_{\text{split}}$, $n_{\text{leaf}}$, $d_{\text{max}}$ & Non-parametric partition; variance reduction. \\
XGBoost & Gradient Boosting & Tabular 2D ($P_H$) & $\lambda_{\text{reg}}$, $\gamma_{\text{split}}$, shrinkage & $\eta$, $d_{\text{max}}$, colsample, $\lambda$, $\alpha_{\text{reg}}$ & Additive boosting; non-linear interaction. \\
LSTM (Linear) & Recurrent Deep & Sequence 3D (10d) & Dropout, weight decay & $d_h, l, dr, bs$ & Decad sequential tracking; linear projection. \\
LSTM (PReLU) & Recurrent Deep & Sequence 3D (10d) & Dropout, weight decay & $d_h, l, dr, bs$ & Decad tracking + temporal attention + PReLU head. \\
\bottomrule
\end{tabularx}
\vspace{0.2em}
{\footnotesize \textit{Source: Author's synthesis of candidate forecasting methodologies.}}
\end{table}

% ==============================================================================
% SECTION 4.4: CHRONOLOGICAL EXPANDING-WINDOW VALIDATION PROTOCOL
% ==============================================================================

\section{Chronological Expanding-Window Validation Protocol}
\label{sec:meth_validation}

\subsection{Three-Phase Temporal Splitting Scheme}
\label{sec:meth_expanding_split}

Random observation splits leak spatial weather signals across contiguous counties \citep{Sweetetal2023}. We enforce a three-phase chronological expanding-window partition (Table~\ref{tab:validation_phases}):
\begin{enumerate}
    \item \textbf{Phase 1: Historical Baseline Training (1951--1984):} 34 years ($4{,}590$ observations) establishing base climatology, secular trends, and initial model weights.
    \item \textbf{Phase 2: Hyperparameter Validation (1985--1995):} 11 years ($1{,}485$ observations) balanced across climatic regimes (6 positive, 5 negative anomaly campaigns), calibrating hyperparameters against severe historical shocks (1988 drought, 1993 flood).
    \item \textbf{Phase 3: Out-of-Sample Expanding-Window Testing (2011--2025):} 15 modern harvest campaigns ($2{,}025$ blind evaluations per horizon). Training expands across $\{1951, \dots, Y-1\}$---incorporating the 1996--2010 historical span into the expanding training pool prior to 2011---to evaluate held-out years 2011--2025 strictly out-of-sample across contemporary extremes (2012 severe drought, 2019 planting crisis, 2023 flash drought) and record bumper harvests (2016, 2021).
\end{enumerate}

\begin{table}[htbp]
\centering
\singlespacing
\footnotesize
\setlength{\tabcolsep}{5pt}
\renewcommand{\arraystretch}{1.1}
\caption{Three-phase chronological validation and expanding-window testing framework.}
\label{tab:validation_phases}
\begin{tabularx}{\textwidth}{l p{7.0cm} X}
\toprule
\textbf{Phase / Fold} & \textbf{Training Partition Span} & \textbf{Target Evaluation Span} \\
\midrule
Phase 1 (Base) & 1951--1984 (34 years, $4{,}590$ county-year obs) & Baseline historical model fitting. \\
Phase 2 (Val)  & 1951--1984 training pool & 1985--1995 (11 years, $1{,}485$ obs; validation and tuning). \\
\midrule
Fold 1 (2011)  & 1951--2010 (60 years, $8{,}100$ county-year obs) & 2011 blind test (135 county evaluations). \\
Fold 2 (2012)  & 1951--2011 (61 years, $8{,}235$ county-year obs) & 2012 blind test (135 obs, Flash Drought). \\
\multicolumn{3}{c}{\dots \dots \dots} \\
Fold 15 (2025) & 1951--2024 (74 years, $9{,}990$ county-year obs) & 2025 blind test (135 county evaluations). \\
\bottomrule
\end{tabularx}
\vspace{0.2em}
{\footnotesize \textit{Source: Author's validation architecture.}}
\end{table}

\subsection{Train-Only Hermetic Isolation}
\label{sec:meth_hermetic_seal}

All transformations compute strictly within each training fold:
\begin{itemize}
    \item Secular trend parameters $(\hat{\alpha}_c^{(Y)}, \hat{\beta}_c^{(Y)})$ are estimated exclusively on historical years $\{1951, \dots, Y-1\}$.
    \item Spatial normalization bounds are extracted solely from training counties.
    \item Weather features are standardized using training partition moments:
    \begin{equation}
        \tilde{w}_{k,j} = \frac{w_{k,j} - \mu_{k,j}^{\text{train}}}{\sigma_{k,j}^{\text{train}}}
        \label{eq:feature_standardization}
    \end{equation}
    Moments $\mathbf{\mu}^{\text{train}}$ and $\mathbf{\sigma}^{\text{train}}$ apply to test year $Y$ without updating.
\end{itemize}

\subsection{Automated Single-Pass Hyperparameter Grid Search}
\label{sec:meth_grid_search}

Hyperparameter tuning executes across validation partition $\mathcal{T}_{\text{val}}$ (1985--1995, $1{,}485$ observations) using chronological expanding-window grid search. For each candidate configuration and validation campaign $Y_{\text{val}} \in \{1985, \dots, 1995\}$, models train on historical campaigns $\{1951, \dots, Y_{\text{val}}-1\}$ with train-only detrending and scaling, predicting year $Y_{\text{val}}$. Table~\ref{tab:hyperparameter_grids} summarizes the search spaces.

\begin{table}[htbp]
\centering
\singlespacing
\footnotesize
\setlength{\tabcolsep}{6pt}
\renewcommand{\arraystretch}{1.1}
\caption{Pre-specified hyperparameter search spaces for expanding-window validation grid search.}
\label{tab:hyperparameter_grids}
\begin{tabularx}{\textwidth}{l p{4.0cm} X c}
\toprule
\textbf{Model} & \textbf{Hyperparameter} & \textbf{Search Space} & \textbf{Combinations} \\
\midrule
\multirow{2}{*}{\textbf{ElasticNet}} & Regularization ($\lambda$) & $\{10^{-3}, \, 10^{-2}, \, 10^{-1}, \, 1.0, \, 5.0\}$ & \multirow{2}{*}{25} \\
 & Penalty Ratio ($\alpha$) & $\{0.1, \, 0.3, \, 0.5, \, 0.7, \, 0.9\}$ & \\
\midrule
\multirow{3}{*}{\textbf{Random Forest}} & Split Features ($p_{\text{split}}$) & $\{\text{'sqrt'}, \, 0.33, \, 0.5\} \times P_H$ & \multirow{3}{*}{18} \\
 & Leaf Samples ($n_{\text{leaf}}$) & $\{5, \, 20\}$ & \\
 & Maximum Depth ($d_{\text{max}}$) & $\{8, \, 14, \, \text{None}\}$ & \\
\midrule
\multirow{4}{*}{\textbf{XGBoost}} & Learning Rate ($\eta$) & $\{0.03, \, 0.06\}$ & \multirow{4}{*}{16} \\
 & Maximum Depth ($d_{\text{max}}$) & $\{3, \, 5\}$ & \\
 & $L_2$ Leaf Weight ($\lambda_{\text{reg}}$) & $\{1.0, \, 5.0\}$ & \\
 & $L_1$ Leaf Penalty ($\alpha_{\text{reg}}$) & $\{0.0, \, 0.5\}$ & \\
\midrule
\multirow{4}{*}{\textbf{LSTM}} & Hidden Dimension ($d_h$) & $\{64, \, 128, \, 256\}$ & \multirow{4}{*}{44} \\
 & Recurrent Layers ($l$) & $\{1, \, 2\}$ & \\
 & Dropout Rate ($dr$) & $\{0.0, \, 0.2, \, 0.4\}$ & \\
 & Batch Size ($bs$) & $\{25, \, 64\}$ & \\
\bottomrule
\end{tabularx}
\vspace{0.2em}
{\footnotesize \textit{Source: Author's specification following domain standards \citep{Torsonietal2023, Sharmaetal2025, Geron2022}}. In Random Forest, $B=100$ trees are used during validation tuning and $B=300$ during final test evaluation. In XGBoost, feature sampling is fixed at $0.7 \times P_H$, row subsampling at $0.8$, and leaf weight threshold at $\text{min\_child\_weight}=3$. In LSTM, learning rate is self-calibrated via dynamic plateau scheduling and early stopping on inner validation folds.}
\end{table}

The optimal parameter vector $\mathbf{\Theta}_H^*$ for horizon $H$ minimizes pooled validation RMSE across the 11 expanding validation folds:
\begin{equation}
    \mathbf{\Theta}_H^* = \arg\min_{\mathbf{\Theta} \in \mathcal{G}} \sqrt{\frac{1}{C \cdot |\mathcal{T}_{\text{val}}|} \sum_{Y_{\text{val}}=1985}^{1995} \sum_{c=1}^C \left( \hat{\epsilon}_{c,Y_{\text{val}}}^{(H)}(\mathbf{\Theta}) - \epsilon_{c,Y_{\text{val}}} \right)^2}
    \label{eq:inner_val_selection}
\end{equation}
Selected hyperparameters remain fixed during the subsequent out-of-sample expanding-window evaluation across 2011--2025, where models re-estimate weights over $\{1951, \dots, Y-1\}$ prior to predicting test year $Y$.

% ==============================================================================
% SECTION 4.5: PERFORMANCE EVALUATION AND OUT-OF-SAMPLE TAIL-RISK STRESS TESTING
% ==============================================================================

\section{Performance Evaluation and Out-of-Sample Tail-Risk Stress Testing}
\label{sec:meth_evaluation}

\subsection{Continuous Loss Metrics and Aggregation Hierarchy}
\label{sec:meth_loss_metrics}

Out-of-sample performance is evaluated on continuous yield anomaly prediction error:
\begin{equation}
    e_{c,Y}^{(H)} = \hat{\epsilon}_{c,Y}^{(H)} - \epsilon_{c,Y} = \hat{y}_{c,Y}^{(H)} - y_{c,Y}
    \label{eq:anomaly_error_equality}
\end{equation}
Because trend extrapolation $\hat{\tau}_{c,Y}$ is identical across models, evaluating errors on $\epsilon_{c,Y}$ is mathematically equivalent to evaluating errors on total yield levels $y_{c,Y}$, without secular trend variance inflating reported accuracy.

Predictions are generated for each individual county $c$ and year $Y$. To synthesize performance across the panel, we evaluate continuous loss metrics over all $N_{\text{test}} = 15$ test years ($M_{\text{test}} = 15 \times 135 = 2{,}025$ evaluations per horizon):
\begin{align}
    \text{RMSE}_{\text{OOS}}(H) &= \sqrt{\frac{1}{M_{\text{test}}} \sum_{Y=2011}^{2025} \sum_{c=1}^C \left( \hat{\epsilon}_{c,Y}^{(H)} - \epsilon_{c,Y} \right)^2} \label{eq:rmse_oos} \\[1ex]
    \text{MAE}_{\text{OOS}}(H) &= \frac{1}{M_{\text{test}}} \sum_{Y=2011}^{2025} \sum_{c=1}^C \left| \hat{\epsilon}_{c,Y}^{(H)} - \epsilon_{c,Y} \right| \label{eq:mae_oos} \\[1ex]
    R^2_{\text{OOS}}(H) &= 1 - \frac{\sum_{Y=2011}^{2025} \sum_{c=1}^C \left( \epsilon_{c,Y} - \hat{\epsilon}_{c,Y}^{(H)} \right)^2}{\sum_{Y=2011}^{2025} \sum_{c=1}^C \epsilon_{c,Y}^2} \label{eq:r2_oos} \\[1ex]
    \text{Skill}_{\text{RMSE}}(H) &= 1 - \frac{\text{RMSE}_{\text{model}}(H)}{\text{RMSE}_{\text{naive}}} \label{eq:skill_score}
\end{align}
$\text{RMSE}_{\text{OOS}}$ serves as the primary metric, penalizing large errors quadratically. $\text{MAE}_{\text{OOS}}$ captures linear error magnitude. $R^2_{\text{OOS}}$ measures explained climatic variance beyond the naive trend: $R^2_{\text{OOS}} = 0$ corresponds to Month 1 naive extrapolation, $R^2_{\text{OOS}} > 0$ establishes genuine meteorological skill, and $R^2_{\text{OOS}} < 0$ reveals overfitting degradation. Skill score $\text{Skill}_{\text{RMSE}}$ tracks percentage error reduction relative to trend extrapolation.

To assess regional stability without presenting 135 separate county tables, we compute cross-sectional annual averages:
\begin{equation}
    \overline{\text{RMSE}}_Y(H) = \sqrt{\frac{1}{C} \sum_{c=1}^C \left( \hat{\epsilon}_{c,Y}^{(H)} - \epsilon_{c,Y} \right)^2}, \quad
    \text{Bias}_Y(H) = \frac{1}{C} \sum_{c=1}^C \left( \hat{\epsilon}_{c,Y}^{(H)} - \epsilon_{c,Y} \right)
    \label{eq:annual_regional_metrics}
\end{equation}
tracking how predictive accuracy and directional bias evolve during extreme shortfall and bumper seasons.

\subsection{Out-of-Sample Tail-Risk Stress Testing}
\label{sec:meth_tail_risk}

To evaluate model resilience during extreme disasters without arbitrary discretization bins, predictions are stress-tested across benchmark historical shock cohorts in the 2011--2025 test window:
\begin{itemize}
    \item \textbf{Adverse Shock Cohorts (Shortfall Campaigns):} The 2012 midsummer flash drought ($-11.7\%$ regional anomaly), the 2019 widespread planting crisis, and the 2023 late-summer flash drought.
    \item \textbf{Favorable Shock Cohorts (Bumper Campaigns):} The modern bumper harvests of 2016 ($+10.8\%$) and 2021 ($+9.8\%$).
\end{itemize}
This protocol evaluates whether flexible regression estimators capture physical shock amplitudes or exhibit conditional shrinkage toward the historical mean.

\subsection{Statistical Inference: Spatial Block-Bootstrap and Diebold-Mariano Test}
\label{sec:meth_significance_bootstrap}

Because synoptic weather systems generate spatial autocorrelation across contiguous Midwestern counties, standard parametric tests produce downward-biased standard errors \citep{Sweetetal2023}. We implement a complete-year spatial block-bootstrap. Resampling operates across complete harvest crop years $Y \in \{2011, \dots, 2025\}$ with replacement ($B_{\text{boot}} = 1{,}000$ iterations). In each bootstrap iteration, all $C = 135$ counties for selected year $Y$ are sampled together, preserving cross-county spatial covariance to generate non-parametric confidence intervals for $\text{RMSE}_{\text{OOS}}(H)$ and $R^2_{\text{OOS}}(H)$.

To test pairwise predictive superiority between competing models and the naive secular benchmark, we compute Diebold-Mariano test statistics on cross-sectionally averaged annual loss differentials:
\begin{equation}
    d_Y(H) = \frac{1}{C} \sum_{c=1}^C \left[ \left(e_{c,Y}^{\text{baseline}}(H)\right)^2 - \left(e_{c,Y}^{\text{model}}(H)\right)^2 \right]
    \label{eq:dm_differential}
\end{equation}
Let $\bar{d}(H) = \frac{1}{N_{\text{test}}} \sum_{Y=2011}^{2025} d_Y(H)$ denote the sample mean differential across the $N_{\text{test}} = 15$ test years. The Diebold-Mariano statistic is defined as:
\begin{equation}
    S_{\text{DM}}(H) = \frac{\bar{d}(H)}{\sqrt{\widehat{\text{Var}}(\bar{d}(H))}}
    \label{eq:dm_statistic}
\end{equation}
where $\widehat{\text{Var}}(\bar{d}(H))$ is estimated via the Newey-West HAC variance estimator. The statistic tests the null hypothesis of equal predictive accuracy ($\mathbb{E}[d_Y] = 0$) against the one-sided alternative of superior model precision.
